% Paleta de colores estandarizada de la tesis GNN-HVA.
% Fuente de verdad unica: toda figura (TikZ, pgfplots) y toda tabla que use
% color debe importar este archivo con % Paleta de colores estandarizada de la tesis GNN-HVA.
% Fuente de verdad unica: toda figura (TikZ, pgfplots) y toda tabla que use
% color debe importar este archivo con % Paleta de colores estandarizada de la tesis GNN-HVA.
% Fuente de verdad unica: toda figura (TikZ, pgfplots) y toda tabla que use
% color debe importar este archivo con % Paleta de colores estandarizada de la tesis GNN-HVA.
% Fuente de verdad unica: toda figura (TikZ, pgfplots) y toda tabla que use
% color debe importar este archivo con \input{tesis-figures/palette} y usar
% los nombres semanticos, nunca HEX sueltos.
%
% Requiere \usepackage{xcolor} en el preambulo del documento.
%
% Las cuatro etapas del pipeline forman una progresion neutro -> destacado:
%   1. Definicion del problema   gris          #8A8D91
%   2. Datos de referencia       amarillo suave #F2C94C
%   3. Optimizacion VQE          naranja        #E8833A
%   4. Prediccion y evaluacion   violeta        #7B54B8
%
% Auxiliares para resultados (buen/mal resultado), accesibles en escala de grises.

% Etapas del pipeline
\definecolor{stageProblem}{HTML}{8A8D91}  % gris    -- etapa 1
\definecolor{stageData}{HTML}{F2C94C}     % amarillo -- etapa 2
\definecolor{stageVQE}{HTML}{E8833A}      % naranja -- etapa 3
\definecolor{stagePred}{HTML}{7B54B8}     % violeta -- etapa 4

% Auxiliares de resultado
\definecolor{resultGood}{HTML}{2E86AB}    % azul    -- configuracion que aprueba
\definecolor{resultBad}{HTML}{C0392B}     % rojo    -- configuracion que falla

% Tintes claros para rellenos de fondo (borde en el color pleno, relleno tenue)
\definecolor{stageProblemBg}{HTML}{E8E9EA}
\definecolor{stageDataBg}{HTML}{FBEFC7}
\definecolor{stageVQEBg}{HTML}{F9DDC7}
\definecolor{stagePredBg}{HTML}{E4DAF2}
 y usar
% los nombres semanticos, nunca HEX sueltos.
%
% Requiere \usepackage{xcolor} en el preambulo del documento.
%
% Las cuatro etapas del pipeline forman una progresion neutro -> destacado:
%   1. Definicion del problema   gris          #8A8D91
%   2. Datos de referencia       amarillo suave #F2C94C
%   3. Optimizacion VQE          naranja        #E8833A
%   4. Prediccion y evaluacion   violeta        #7B54B8
%
% Auxiliares para resultados (buen/mal resultado), accesibles en escala de grises.

% Etapas del pipeline
\definecolor{stageProblem}{HTML}{8A8D91}  % gris    -- etapa 1
\definecolor{stageData}{HTML}{F2C94C}     % amarillo -- etapa 2
\definecolor{stageVQE}{HTML}{E8833A}      % naranja -- etapa 3
\definecolor{stagePred}{HTML}{7B54B8}     % violeta -- etapa 4

% Auxiliares de resultado
\definecolor{resultGood}{HTML}{2E86AB}    % azul    -- configuracion que aprueba
\definecolor{resultBad}{HTML}{C0392B}     % rojo    -- configuracion que falla

% Tintes claros para rellenos de fondo (borde en el color pleno, relleno tenue)
\definecolor{stageProblemBg}{HTML}{E8E9EA}
\definecolor{stageDataBg}{HTML}{FBEFC7}
\definecolor{stageVQEBg}{HTML}{F9DDC7}
\definecolor{stagePredBg}{HTML}{E4DAF2}
 y usar
% los nombres semanticos, nunca HEX sueltos.
%
% Requiere \usepackage{xcolor} en el preambulo del documento.
%
% Las cuatro etapas del pipeline forman una progresion neutro -> destacado:
%   1. Definicion del problema   gris          #8A8D91
%   2. Datos de referencia       amarillo suave #F2C94C
%   3. Optimizacion VQE          naranja        #E8833A
%   4. Prediccion y evaluacion   violeta        #7B54B8
%
% Auxiliares para resultados (buen/mal resultado), accesibles en escala de grises.

% Etapas del pipeline
\definecolor{stageProblem}{HTML}{8A8D91}  % gris    -- etapa 1
\definecolor{stageData}{HTML}{F2C94C}     % amarillo -- etapa 2
\definecolor{stageVQE}{HTML}{E8833A}      % naranja -- etapa 3
\definecolor{stagePred}{HTML}{7B54B8}     % violeta -- etapa 4

% Auxiliares de resultado
\definecolor{resultGood}{HTML}{2E86AB}    % azul    -- configuracion que aprueba
\definecolor{resultBad}{HTML}{C0392B}     % rojo    -- configuracion que falla

% Tintes claros para rellenos de fondo (borde en el color pleno, relleno tenue)
\definecolor{stageProblemBg}{HTML}{E8E9EA}
\definecolor{stageDataBg}{HTML}{FBEFC7}
\definecolor{stageVQEBg}{HTML}{F9DDC7}
\definecolor{stagePredBg}{HTML}{E4DAF2}
 y usar
% los nombres semanticos, nunca HEX sueltos.
%
% Requiere \usepackage{xcolor} en el preambulo del documento.
%
% Las cuatro etapas del pipeline forman una progresion neutro -> destacado:
%   1. Definicion del problema   gris          #8A8D91
%   2. Datos de referencia       amarillo suave #F2C94C
%   3. Optimizacion VQE          naranja        #E8833A
%   4. Prediccion y evaluacion   violeta        #7B54B8
%
% Auxiliares para resultados (buen/mal resultado), accesibles en escala de grises.

% Etapas del pipeline
\definecolor{stageProblem}{HTML}{8A8D91}  % gris    -- etapa 1
\definecolor{stageData}{HTML}{F2C94C}     % amarillo -- etapa 2
\definecolor{stageVQE}{HTML}{E8833A}      % naranja -- etapa 3
\definecolor{stagePred}{HTML}{7B54B8}     % violeta -- etapa 4

% Auxiliares de resultado
\definecolor{resultGood}{HTML}{2E86AB}    % azul    -- configuracion que aprueba
\definecolor{resultBad}{HTML}{C0392B}     % rojo    -- configuracion que falla

% Tintes claros para rellenos de fondo (borde en el color pleno, relleno tenue)
\definecolor{stageProblemBg}{HTML}{E8E9EA}
\definecolor{stageDataBg}{HTML}{FBEFC7}
\definecolor{stageVQEBg}{HTML}{F9DDC7}
\definecolor{stagePredBg}{HTML}{E4DAF2}
